\section{Experiments}
This section reports the results that are actually supported by the current repository release. We evaluate synthetic-only training and synthetic-plus-weaklabel training. Real-pilot and raw-video experiments remain out of scope for this version.

\subsection{Metrics}
Timing prediction is evaluated with accuracy, balanced accuracy, macro-F1, and per-class precision/recall/F1. Because the timing labels are imbalanced, balanced accuracy and macro-F1 are the primary metrics. Strategy prediction is evaluated with strategy accuracy and strategy macro-F1. For the multitask model, we additionally report state MAE. For the tri-modal emotion encoder, we report fine-emotion accuracy, balanced accuracy, macro-F1, valence MAE, arousal MAE, and auxiliary state MAE.

\subsection{Protocol}
The formal experiments reported here were trained on two NVIDIA A800 80GB GPUs using the \texttt{remote\_a800.yaml} setting. Synthetic training uses a 24{,}000/2{,}400/2{,}400 split; mixed training augments the synthetic train split with 286 weak-label training episodes extracted from historical logs and evaluates on both the held-out synthetic split and the held-out weaklabel split. The bridge-log parser is implemented but not included in the main experiments because its current volume is too small for stable training conclusions.

Two caveats are essential for interpreting the tables. First, the synthetic benchmark is now a real three-class task, but it is still rule-generated and therefore measures rule-recovery plus temporal generalization rather than natural in-the-wild behavior. Second, the combined test columns in Table~\ref{tab:main-results} merge the 2{,}400-window synthetic test set with the 62-window weaklabel test set, so they are dominated by synthetic windows and should be treated as aggregate diagnostics rather than as a proxy for real-world robustness.

\subsection{Main Results}
\begin{table*}[t]
\centering
\caption{Main offline results. Combined columns use the full held-out test split produced by each training setting; because that split contains 2{,}400 synthetic windows but only 62 weaklabel windows, the combined columns are synthetic-dominated. Synthetic columns evaluate only the held-out synthetic split, and the weaklabel column reports timing Macro-F1 on the held-out weaklabel split when available. Strategy scores remain template-like and should be interpreted as constrained planning sanity checks rather than open-ended response quality. All numbers are percentages except state MAE.}
\label{tab:main-results}
\small
\begin{tabular}{llccccccc}
\toprule
Model & Train Data & Comb. Bal. Acc. & Comb. Macro-F1 & Synth Bal. Acc. & Synth Macro-F1 & Weaklabel Macro-F1 & Comb. Strategy Macro-F1 & State MAE \\
\midrule
Structured & Synthetic & 85.56 & 79.11 & 85.56 & 79.11 & -- & 80.44 & -- \\
Structured & Synthetic + weaklabel & 83.81 & 77.78 & 84.98 & 78.72 & 27.29 & 78.99 & -- \\
Multitask temporal & Synthetic & 86.88 & 77.21 & 86.88 & 77.21 & -- & -- & 0.069 \\
Multitask temporal & Synthetic + weaklabel & 82.24 & 76.15 & 86.66 & 78.17 & 7.62 & -- & 0.083 \\
Joint temporal & Synthetic + weaklabel & 85.79 & 80.45 & 86.90 & 81.06 & 36.29 & 82.66 & -- \\
\bottomrule
\end{tabular}
\end{table*}


Three findings are robust across the current runs. First, on the held-out synthetic split the temporal models are now credible three-class timing baselines, and the joint branch is strongest overall at 86.90\% balanced accuracy and 81.06\% timing Macro-F1, with the structured sanity baseline close behind at 85.56\% and 79.11\%. Second, mixed training improves the aggregate offline numbers for the joint branch, but the combined columns remain synthetic-heavy; the strongest held-out weaklabel timing Macro-F1 is still only 36.29\%. Third, the multitask branch remains important because it is the only current model that predicts latent state directly, reaching a state MAE of 0.069 on synthetic-only training and 0.083 under mixed training.

\subsection{Weaklabel Generalization}
\begin{table}[t]
\centering
\caption{Held-out weaklabel results on the 62-window database-derived split extracted from historical logs.}
\label{tab:weaklabel-results}
\small
\begin{tabular}{lccc}
\toprule
Model & Bal. Acc. & Macro-F1 & Strategy Macro-F1 \\
\midrule
Structured & 32.05 & 27.29 & 22.50 \\
Multitask temporal & 33.33 & 7.62 & -- \\
Joint temporal & 58.67 & 36.29 & 18.06 \\
\bottomrule
\end{tabular}
\end{table}


Held-out weaklabel performance remains the hardest setting. The exploratory joint model is strongest on weaklabel timing Macro-F1 at 36.29\%, the structured baseline reaches 27.29\%, and the multitask model drops to 7.62\%. Strategy transfer is weaker still: the best held-out weaklabel strategy Macro-F1 is only 22.50\% from the structured baseline, while the joint model reaches 18.06\%. These numbers are useful as a bootstrapping signal, not as proof of real-world robustness. They mainly quantify the current synthetic-to-log mismatch and motivate future raw-video ingestion, stronger weaklabel reconstruction, and real-pilot collection.

\subsection{Tri-Modal Emotion Benchmark}
\begin{table}[t]
\centering
\caption{Synthetic tri-modal fine-grained emotion benchmark. Full fusion uses face, motion, and audio together; ablations zero out the other two modalities at test time. All accuracy and F1 numbers are percentages.}
\label{tab:trimodal-results}
\small
\begin{tabular}{lccccc}
\toprule
Model & Acc. & Macro-F1 & Val. MAE & Aro. MAE & State MAE \\
\midrule
Full tri-modal & 97.59 & 97.58 & 0.068 & 0.055 & 0.053 \\
Face only & 24.33 & 11.08 & 0.209 & 0.452 & 0.181 \\
Motion only & 8.33 & 1.28 & 0.302 & 0.478 & 0.242 \\
Audio only & 31.11 & 16.65 & 0.302 & 0.150 & 0.244 \\
\bottomrule
\end{tabular}
\end{table}


The new tri-modal benchmark tests the perception layer in isolation. Full fusion reaches 97.59\% accuracy and 97.58\% Macro-F1 on the held-out synthetic benchmark, while also keeping valence MAE at 0.068 and arousal MAE at 0.055. The ablations are more informative than the headline number: face-only, motion-only, and audio-only settings collapse to 24.33\%, 8.33\%, and 31.11\% accuracy, respectively. On this controlled benchmark, fusion is not a marginal gain but the core source of performance.

\subsection{Tri-Modal Care Relation Benchmark}
\begin{table}[t]
\centering
\caption{Synthetic tri-modal care-need and care-timing benchmark. Full fusion uses face, motion, and audio together; ablations zero out the other two modalities at test time. All classification numbers are percentages.}
\label{tab:trimodal-care-results}
\small
\begin{tabular}{lcccc}
\toprule
Model & Need AUROC & Need F1 & Timing Macro-F1 & Support MAE \\
\midrule
Full tri-modal & 99.45 & 97.85 & 94.00 & 0.058 \\
Face only & 98.83 & 96.82 & 51.88 & 0.073 \\
Motion only & 51.89 & 0.00 & 16.90 & 0.374 \\
Audio only & 79.84 & 0.00 & 16.90 & 0.350 \\
\bottomrule
\end{tabular}
\end{table}


The second tri-modal benchmark asks a harder question: can the same multimodal encoder predict whether the user needs care and whether the appropriate action is \texttt{none}, \texttt{delay}, or \texttt{immediate}? Full fusion reaches 99.45\% care-need AUROC, 97.85\% care-need F1, and 94.00\% care-timing Macro-F1. The ablations reveal a more nuanced picture than in the emotion benchmark. Face-only remains surprisingly strong for binary care-need prediction, but its timing Macro-F1 drops to 51.88\%, while motion-only and audio-only collapse almost entirely. On this synthetic benchmark, binary concern detection is easier than timing-specific intervention planning, and full fusion matters more for the latter.

\subsection{Status Audit}
\begin{table*}[t]
\centering
\caption{Status audit for the current release. This table separates validated components from pending ones so the paper can state precisely what has been established and what remains incomplete.}
\label{tab:audit}
\small
\begin{tabular}{p{0.24\linewidth}p{0.16\linewidth}p{0.54\linewidth}}
\toprule
Component & Status & Evidence / Interpretation \\
\midrule
Synthetic benchmark & Complete & Train/dev/test JSONL datasets and reproducible generators are in the repository and are used in all reported experiments. \\
Weak-label extraction & Complete & \texttt{auth.db} logs are converted into 286/61/62 weak-label train/dev/test episodes and used in the mixed experiments. \\
Structured baseline & Complete & Trained and evaluated; it remains a strong sanity baseline and a useful transfer comparison against the temporal models. \\
Temporal multitask model & Complete & Trained on both synthetic-only and synthetic+weaklabel regimes; this is the main state-prediction branch and a competitive neural timing baseline. \\
Joint timing-to-strategy model & Exploratory but reported & Trained and evaluated; it is currently the strongest offline timing branch, but weaklabel strategy generalization remains low. \\
Raw video backbone & Pending & The current release still uses temporal descriptors rather than raw frame or waveform encoders. \\
Real pilot / human study & Pending & No human evaluation results are reported yet, so claims must remain offline and non-clinical. \\
Top-line deployment claim & Not allowed & Current evidence does not support claims about real-user efficacy, cross-scene robustness, or therapeutic benefit. \\
\bottomrule
\end{tabular}
\end{table*}

